% =====================================================================
% SEC 05 — Discussão
% =====================================================================
\section{Discussão}
\label{sec:discussao}

\subsection{Por que o esquema posicional empata com o \emph{top-k}}

O resultado central — Recamán $=$ \emph{top-k} em diversidade — tem uma causa
\emph{estrutural}, não acidental. O diversificador pós-Recamán (Algoritmo
\ref{alg:recaman}) opera \emph{sobre posições} do ranking: ele salta para índices
determinísticos $\mathrm{offset}_m = (1+a_m)\bmod N$, mas \textbf{não consulta o
conteúdo} (âncora/fonte) dos itens. Quando o ranking é \emph{localmente
monopolista} — i.e., quando várias posições consecutivas são ocupadas por
documentos da \emph{mesma família de perspectiva} — os offsets que caem nessas
posições selecionam novamente a mesma família.

No corpus-piloto, as três primeiras posições de cada ranking são preenchidas por
três documentos da mesma família (ex.: \texttt{family\_2}). Os offsets iniciais
de Recamán e o fato de o algoritmo partir do top-1 fazem com que a seleção recaia
majoritariamente nessa família, reproduzindo o \emph{top-k}. Formalmente, um
diversificador posicional $P$ só garante diversidade de $\Div$ quando o mapeamento
posição$\to$família é ``colorida'', o que não ocorre em rankings monopolistas.

\subsection{Contraste com o MMR}

O MMR, por maximizar explicitamente
$\lambda\,\mathrm{relev} - (1-\lambda)\max\Sim(\cdot,\cdot)$, usa \emph{dissimilaridade de
conteúdo} entre os itens já selecionados e os candidatos. Isso o torna capaz de
\emph{quebrar} o monopolismo: mesmo que a família dominante ocupe o topo, o MMR
prefere um candidato de outra família desde que a dissimilaridade compense a perda
de relevância. O custo disso é uma \emph{queda de relevância} ($\approx 6{,}7\%$ no
corpus-piloto), acima da tolerância definida neste estudo — o que evidencia o
`trade-off` clássico entre diversidade e relevância.

\subsection{Implicações para diversificação em RAG}

Do ponto de vista prático, o achado sugere:

\begin{itemize}
  \item Diversificadores \emph{posicionais} (como a variante pós-Recamán aqui
        avaliada) são eficazes para \emph{espalhar por posição}, mas não garantem
        diversidade de \emph{fonte/âncora} quando o ranqueador upstream concentra
        famílias no topo. Uma direção de melhoria é \emph{operar o diversificador
        sobre âncoras}, recomputando os offsets no espaço de famílias (não de
        posições) — além do escopo deste artigo.
  \item Para consultas de paisagem/multicunha em RAG, uma combinação
        \emph{relevância+dissimilaridade} (na linha do MMR) tende a ser mais
        adequada, aceitando-se uma queda controlada de relevância, desde que a
        métrica de diversidade seja monitorada explicitamente (Eq.
        \ref{eq:diversity}).
  \item A métrica $\Div(\mathcal{S})$ e o \emph{benchmark} propostos provêm uma
        fundação para \emph{medir} — e não apenas \emph{assumir} — o ganho de
        diversificação em futuros trabalhos.
\end{itemize}

\subsection{Limitações do estudo}

Declaramos as limitações de forma explícita (anti-overclaim):

\begin{enumerate}
  \item \textbf{Corpus-piloto sintético:} os resultados referem-se a um corpus
        controlado e determinístico; \textbf{não generalizam} para produção nem
        para corpora reais de grande escala.
  \item \textbf{Sem embeddings externos:} o ranqueador e a similaridade do MMR
        usam sobreposição de tokens, não representações densas; comportamento com
        embeddings pode diferir.
  \item \textbf{Sem geração de linguagem:} avaliamos \emph{seleção/recuperação},
        não a qualidade do texto gerado pelo LLM \emph{a partir} das evidências.
  \item \textbf{Conjunto de queries pequeno} ($n=4$): os agregados têm caráter
        ilustrativo; um experimento com dezenas/centenas de queries e múltiplas
        sementes é necessário para inferência estatística.
  \item \textbf{Parâmetro $\lambda=0{,}7$ fixo} no MMR; a sensibilidade a
        $\lambda$ não foi explorada aqui.
\end{enumerate}

\subsection{Trabalho futuro}

Como extensões naturais, propomos: (i) o redesenho do diversificador para operar
sobre \emph{âncoras canônicas} (offsets no espaço de famílias); (ii) a validação
em corpus real (ex.: coleções científicas abertas), com mais queries e sementes;
(iii) a inclusão de embeddings densos e similaridade semântica; e (iv) a
avaliação \emph{extrínseca} da qualidade da resposta gerada sob as diferentes
estratégias de seleção.
